% ============================================================================
\clearpage
\section*{v4.10.0-r36 Source-Only: A Vacuum Background as Carrier of the Directed OP5 Mediator}
\addcontentsline{toc}{section}{v4.10.0-r36 Source-Only: A Vacuum Background for the OP5 Mediator}
\label{sec:v4100-r36-background-mediator}
% ============================================================================

\paragraph{Purpose.}
After v4.10.0-r35 the material route carries the directed contrast of the OP5 case only partially: a density coupling gives the numerator
\(a_\perp\) instead of \(\tfrac23a_\perp^2\).  This layer follows the author's choice of the vacuum route: a new field \(X\) coupled linearly and gauge
invariantly to \(|\Phi|^2\), hence to the Higgs radial mode \(\eta\) via \(|\Phi|^2=v^2+2v\eta+\dots\), in a vacuum that may carry a frame \(u\) and an axis
\(n\).  It is an exploration with a method fixed before the computation and is not blind.  Earlier layers are not rewritten.

\noindent\textbf{Source-only status.}
This v4.10.0-r36 layer is source-only.  It does not build a PDF, does not assign a new DOI,
and does not authorize a public release.  The v5.3.0 Version DOI
\href{https://doi.org/10.5281/zenodo.24434403}{10.5281/zenodo.24434403} and the Concept DOI
\href{https://doi.org/10.5281/zenodo.19210728}{10.5281/zenodo.19210728} are unchanged by it.

\subsection*{v4.10.0-r36 results}
\label{subsec:v4100-r36-results}

Signature \((+,-,-,-)\), \(p=(\omega,\mathbf k)\), \(x=\omega^2\), \(a_\perp=k_1^2+k_2^2\), \(u=(1,0,0,0)\), \(n=(0,0,0,1)\), \(T^{\mu\nu}=u^\mu u^\nu\),
\(S^{\mu\nu}=\operatorname{diag}(0,1,1,0)\).  The numerator of the mediator term is the contraction of the vertex \(1\), \(p^\mu\) or \(p^\mu p^\nu\) with the
polarization sum.

\paragraph{Theoremlet v4.10.0-r36.1 (no directed pole in a Lorentz-invariant vacuum).}
For \(X|\Phi|^2\) the pole has a constant numerator.  For \(X_\mu\partial^\mu|\Phi|^2\) with the Proca sum \(P^{\mu\nu}=-\eta^{\mu\nu}+p^\mu p^\nu/M^2\) the numerator is
\(p^2(p^2-M^2)/M^2\), and the exchange is the contact term \(p^2/M^2\).  For \(X_{\mu\nu}\partial^\mu\partial^\nu|\Phi|^2\) with the Fierz--Pauli sum the numerator is
\(\tfrac23(p^2)^2(p^2-M^2)^2/M^4\), with a double zero at the pole.  Hence no linear coupling of a new field to \(|\Phi|^2\) gives a directed pole in a
Lorentz-invariant vacuum.  \emph{Proof.} \(p_\mu P^{\mu\nu}=p^\nu(p^2-M^2)/M^2\); insert into the polarization sums.  \(\square\)

\paragraph{Theoremlet v4.10.0-r36.2 (background building blocks).}
With polarizations restricted by \(T\) and \(S\), a vector transverse to the axis gives the numerator \(a_\perp\).  For a symmetric tensor the blocks
\(TT\), the time--transverse block \(B\), the transverse traceless block \(G\), the transverse trace block \(D\) and the mixing \(E=TS+ST\) contribute
\(x^2\), \(2a_\perp x\), \(a_\perp^2/2\), \(a_\perp^2/2\) and \(2a_\perp x\).  \emph{Proof.} Direct contraction with \(p_\mu p_\nu p_\rho p_\sigma\).  \(\square\)

\paragraph{Theoremlet v4.10.0-r36.3 (the numerator \(N_m\) with a positive polarization sum).}
The sum \(\alpha\,TT+\beta B+\gamma G+\delta D+\varepsilon E\) gives \(N_m=\tfrac23(x^2-a_\perp x+a_\perp^2)\) exactly if
\[
 \alpha=\tfrac23,\qquad \gamma=\tfrac{4-3\delta}{3},\qquad \varepsilon=-\tfrac{1+3\beta}{3},
\]
and it is positive semidefinite exactly for \(0\le\beta<\tfrac13\), \((1+3\beta)^2/3\le\delta\le\tfrac43\), or \(\beta=\tfrac13\), \(\delta=\tfrac43\).  At that end point
\(\gamma=0\), \(\varepsilon=-\tfrac23\), and the block of \(TT\) and transverse trace has rank one.  The sum with the metric of the subspace orthogonal to \(n\)
gives \((1-c)(x-a_\perp)^2\) instead, and \(N_m-\tfrac23(x-a_\perp)^2=\tfrac23a_\perp x\).
\emph{Proof.} Coefficient comparison in \(x^2,a_\perp x,a_\perp^2\); positivity of the orthogonal blocks \(B,G\) and of the \(2\times2\) block
\(\bigl(\begin{smallmatrix}\alpha&\sqrt2\varepsilon\\\sqrt2\varepsilon&\delta\end{smallmatrix}\bigr)\), solved exactly.  \(\square\)

\paragraph{Remark v4.10.0-r36.4 (reading).}
Against the criteria fixed beforehand the vacuum route carries the directed contrast \emph{partially}.  A directed pole needs a background in
every case (r36.1).  With frame and axis the numerator \(N_m\) is reachable exactly, including the power \(a_\perp^2\) that the material route lacks
(r35.3); its negative mixed term comes from a coupling between the time--time component and the transverse trace of the field.  Set rather
than derived are the frame and the axis, the polarization structure (a two-parameter family) and an action producing it.  Candidates for the
background named by the Companion are the flow and the spin vector of the OP5 seed (spin along the azimuthal direction) and the spacelike
tachyon image of a matter state at rest; both are readings.

\subsection*{v4.10.0-r36 frames}
\label{subsec:v4100-r36-frames}

All quantities are given in the matter frame M, the background vectors included.  No quantity of another frame is compared.

\subsection*{v4.10.0-r36 register entries}
\label{subsec:v4100-r36-register}

\begin{center}
\small
\begin{tabular}{@{}>{\raggedright\arraybackslash}p{0.34\textwidth}>{\raggedright\arraybackslash}p{0.58\textwidth}@{}}
\toprule
statement & status after v4.10.0-r36\\
\midrule
directed pole without background & excluded for linear couplings to \(|\Phi|^2\) (r36.1)\\
\(N_m\) from a vacuum background & reachable with a positive spin-2 sum; structure, frame and axis set (r36.3)\\
origin of the background & open; OP5 seed and tachyon image as readings\\
\bottomrule
\end{tabular}
\end{center}

\subsection*{v4.10.0-r36 notes}
\label{subsec:v4100-r36-notes}

\texttt{notes/\allowbreak research/\allowbreak op5\_\allowbreak background\_\allowbreak mediator\_\allowbreak 2026\_\allowbreak 11\_\allowbreak 22/} (method note commit 9f91c274,
script \texttt{background\_\allowbreak mediator.wl}, output in \texttt{run/}).

\subsection*{v4.10.0-r36 ledger and guard update}
\label{subsec:v4100-r36-ledger-guard}

\[
R_{\rm chk}^{v4100\text{-}r36}=(1,1,0)
\]
for (status recorded; computations checked in the cited script; the background and an action for the polarization sum not derived).

The layer spends no guard:
\[
R_{\rm glob}^{(v4100\text{-}r36,+)}=(0,1,1,1,1),
\qquad
\Gamma_{\rm guard}^{v4100\text{-}r36}=(1,1,1,1,1),
\qquad
\beta_4=1/256\ \text{unspent}.
\]
The layer does not imply exact finite \(C5\) zero closure, exact all-readout terminality, \(C6/T2\), a global photon-clock/RH theorem,
or global all-readout closure.  OP1--OP5, C5 and C6/T2 remain open.
